\chapter{Personal, licencias y formación}\label{ch:19}
\objetivos{Separar conocimiento académico, licencia, habilitación y autorización; interpretar el propósito de las categorías sin simplificar sus privilegios.}
\section{Cuatro elementos diferentes}
Una titulación universitaria acredita una formación académica. Una licencia de mantenimiento acredita una condición regulada bajo un sistema y categoría. Una habilitación o anotación delimita productos o alcance según corresponda. Una autorización de certificación de una organización vincula a la persona con funciones dentro de sus procedimientos. Ninguno de estos elementos sustituye automáticamente a todos los demás.

Parte 66 regula la licencia de mantenimiento de aeronaves dentro de su ámbito. Parte 147 regula organizaciones de formación y exámenes de mantenimiento. La aprobación de una escuela no equivale a una aprobación de taller, y superar un curso no significa que ya se posean todos los privilegios de certificación sobre cualquier aeronave.
\section{Categorías y privilegios}
Las categorías incluyen A, B1, B2, B2L, B3, L y C, con subcategorías, límites y condiciones. En términos orientativos, A se relaciona con determinadas tareas menores de línea y rectificación simple dentro de la autorización; B1 con estructura, planta de potencia y sistemas mecánicos y eléctricos y ciertos trabajos de aviónica limitados; B2 con ámbitos eléctricos y de aviónica y las condiciones especificadas. La letra de la categoría no permite omitir el texto de 66.A.20.

Las categorías B2L, B3 y L tienen poblaciones o sistemas específicos. C se relaciona con certificación tras mantenimiento de base dentro de las condiciones aplicables. Una descripción en una línea sirve de mapa, no de autorización. La licencia puede contener limitaciones y la organización debe evaluar competencia y emitir la autorización que corresponda.
\begin{table}[htbp]\centering\small\caption{Documentos personales y pregunta que resuelven.}\label{tab:personal}
\begin{tabular}{p{.26\linewidth}p{.59\linewidth}}\toprule
Elemento & Pregunta principal\\\midrule
Título académico & ¿Qué formación ha acreditado la institución?\\[4pt]
Licencia Parte 66 & ¿Qué categoría, limitaciones y condiciones reconoce el sistema de licencia?\\[4pt]
Anotación de tipo & ¿Qué tipo o grupo y qué condiciones constan en la licencia?\\[4pt]
Autorización de organización & ¿Qué puede certificar esa persona en la organización, con qué alcance y procedimientos?\\\bottomrule
\end{tabular}\end{table}
\section{Conocimientos y experiencia}
Los puntos 66.A.25 y 66.A.30 tratan conocimientos básicos y experiencia. Las rutas, periodos, créditos y condiciones dependen de la categoría y formación previa. No se da un único número de años como respuesta universal. La experiencia debe corresponder a las tareas y población definidas y poder documentarse de forma aceptable.

El registro debe permitir identificar fecha, producto, tarea, participación y supervisión según la finalidad del documento. Una lista de asistencia a clase no es una experiencia de mantenimiento ejecutado. Un ejercicio de laboratorio universitario puede aportar aprendizaje, pero su reconocimiento reglamentario necesita la vía correspondiente.
\section{Habilitaciones, formación de tipo y continuidad}
66.A.45 regula anotaciones y condiciones relacionadas. La formación teórica y práctica de tipo y, cuando proceda, entrenamiento en el puesto responden a objetivos diferentes. Una evaluación debe comprobar capacidad, no solo asistencia. Las condiciones de ejercicio de privilegios incluyen los requisitos de actualidad, competencia y otros límites del punto aplicable.

La tecnología y los procedimientos cambian; la formación continua se relaciona con productos, procedimientos y factores humanos. Un profesional competente necesita saber detener una actuación fuera de su alcance y buscar datos o apoyo. La capacidad de reconocer incertidumbre forma parte del trabajo técnico.
\section{Qué aporta Parte 147}
Parte 147 establece requisitos de instalaciones, personal, registros, organización de formación y exámenes y procedimientos. Su objeto es dar un marco controlado a la formación aprobada y a las evaluaciones correspondientes. La aprobación tiene alcance, ubicación y limitaciones.

Los exámenes deben preservar integridad y trazabilidad. El contenido de esta obra son ejercicios originales de estudio: no se atribuye equivalencia con bancos oficiales ni con pruebas de habilitación. La colección tampoco entrega un trabajo evaluable para que se presente como realización personal.
\caso{M19. Un grado no crea una firma autorizada}{Una persona acaba el Grado en Ingeniería Aeroespacial y quiere emitir un CRS de una intervención en un motor instalado. Se identifica el sistema regulatorio, tarea, categoría y habilitación, experiencia, competencia y autorización organizativa exigibles. El título puede resultar relevante para una ruta, pero no demuestra que esos elementos estén satisfechos. La conclusión académica es un mapa de requisitos por verificar, no una denegación o concesión administrativa de licencia.}
\section{Cómo estudiar esta parte}
Conviene construir casos en lugar de memorizar únicamente siglas. Para cada persona ficticia se anotan formación, licencia, limitaciones, tipo, experiencia reciente, organización y tarea solicitada. Se comprueba cada condición con su punto y se evita inferir que una condición satisfecha compensa automáticamente otra que falta.

\section{Conocimientos, experiencia, tipo y autorización}
Parte 66 no se reduce a aprobar un examen. La ruta combina conocimientos, experiencia, categoría y, cuando corresponda, habilitación o anotación de tipo y condiciones para ejercer privilegios. 66.A.20 define privilegios y condiciones; .25 y .30, conocimientos y experiencia; .45, habilitaciones. Se lee el escenario concreto sin inventar una cifra de experiencia general para todos.

La autorización organizativa conecta persona y funciones dentro de un alcance. Una licencia no demuestra autorización para cualquier tarea y modelo; trabajar en un taller tampoco concede automáticamente privilegios de certificación. La evaluación de competencia y sus evidencias siguen los requisitos y procedimientos pertinentes.
\section{Formación y equivalencias demostradas}
Parte 147 regula organizaciones de formación en su ámbito. Reconocimiento de curso y certificados tiene condiciones. La universidad aporta bases científicas relevantes, pero el grado no se convierte por nombre en licencia. Posibles créditos o rutas requieren disposición concreta y evaluación competente.

Un expediente separa formación básica, formación de tipo, prácticas, experiencia y autorizaciones. Conserva fechas y contenidos sin inventar equivalencias. El laboratorio puede desarrollar habilidades valiosas con alcance diferente de la experiencia acreditada exigida por una ruta.
\section{Competencia continua y factores humanos}
Cambios de producto, datos y útiles requieren actualización. Familiaridad con una variante puede favorecer suposiciones incorrectas sobre otra. Trabajo reciente y condiciones de ejercicio se consideran en su régimen. Fatiga, comunicación, presión e interrupciones afectan ejecución.

Comprender el mecanismo material no sustituye el procedimiento; conocer una lista sin su propósito también puede generar errores ante una anomalía. La formación conecta conocimiento, ejecución y reconocimiento de límites.
\repaso{¿Carece de utilidad el grado para mantenimiento por no ser licencia?}{La base de análisis y materiales es relevante. Utilidad formativa y privilegio legal son conceptos distintos; se identifican los requisitos y evidencias del paso siguiente.}

\resumen{Título, licencia, habilitación y autorización tienen funciones distintas. Las categorías necesitan leerse con sus privilegios y condiciones. Parte 147 controla formación y examen; Parte 66, licencias dentro de su ámbito. Los casos académicos no crean derechos profesionales.}
\fuentes{Reglamento 1321/2014, 66.A.15, .20, .25, .30 y .45; Parte 147 y 145.A.35 \cite{EU1321}.}
